\section{Frequency cutoffs for all compact smooth tests on generic lines}
\label{sec:regularized-frequency}

The generic-line logarithmic derivative in MA-RS-001 is better
controlled than the reciprocal multiplier itself. This distinction
allows a further explicit regularization: each finite-frequency
expression is an actual trace-class operator, and its trace has a
limit for every compact smooth test. The construction below uses two
bounded compositions before any kernel or trace is evaluated.
Its limit is a limit of traces. It does not assert convergence of
the operators on a merely generic line.

Use the notation of (MA-RS.8), with
$\sigma\in\mathscr G$, $a=\sigma-1/2$, and
\[
 A(t)=\zeta_S(\sigma+it),\qquad
 j(t)=A(t)^{-1},\qquad
 G(z)=\mathcal Mg(z),\qquad m(t)=G(a+it).
 \tag{MA-RS.52}
\]
Here $g\in C_c^\infty(\mathbb R)$ is arbitrary.
Let $\phi$ be bounded and smooth, with
$\phi(-\infty)=1$, $\phi(+\infty)=0$, and
$\phi'\in C_c^\infty(\mathbb R)$, and put
\[
 T_\phi=V_aM_\phi V_a^{-1}
        =\mathcal F_+M_\phi\mathcal F_+^{-1},
 \qquad
 p(\xi)=\int_{\mathbb R}e^{i\xi x}\phi'(x)\,dx.
 \tag{MA-RS.53}
\]
In particular $T_\phi$ is bounded on $L^2(\mathbb R)$; it
is the Fourier conjugate of multiplication by the physical cutoff.
It is not multiplication by $\phi$ in the frequency variable.
Fix a real $\chi\in C_c^\infty(\mathbb R)$ equal to one in a
neighborhood of zero, and set $\chi_N(t)=\chi(t/N)$ for integers
$N\ge1$.

\paragraph{MA-RS-007 (actual finite-frequency traces and their smooth limit).}
Define the operator $K_N(g)$ on $\mathcal H_a$ by the literal
bounded-composition formula
\[
 \begin{split}
 \widetilde K_N(g)&:=V_aK_N(g)V_a^{-1}\\
  &=M_{m\chi_N}\,T_\phi\,M_{\chi_N}
       -M_{mA\chi_N}\,T_\phi\,M_{j\chi_N}.
 \end{split}
 \tag{MA-RS.54}
\]
Both displayed compositions are bounded. Their difference is trace
class, with smooth compactly supported kernel
\[
 \mathcal K_N(t,u)
   =\frac{i}{2\pi}\chi_N(t)\chi_N(u)m(t)A(t)
          \frac{j(t)-j(u)}{t-u}\,p(t-u).
 \tag{MA-RS.55}
\]
The divided difference is assigned its smooth value $j'(t)$
on the diagonal. The trace and its limit are
\[
 \begin{aligned}
 \operatorname{Tr}K_N(g)
   &=-\frac1{2\pi}\int_{\mathbb R}\chi_N(t)^2G(a+it)
                    \frac{\zeta_S'}{\zeta_S}(\sigma+it)\,dt,\\
 \lim_{N\to\infty}\operatorname{Tr}K_N(g)
   &=-\frac1{2\pi}\int_{\mathbb R}G(a+it)
                    \frac{\zeta_S'}{\zeta_S}(\sigma+it)\,dt\\
   &=P_{\overline S}^{+}(g)-G(1/2)
       +\sum_{\operatorname{Re}\rho>\sigma}
                      m_\rho G(\rho-1/2).
 \end{aligned}
 \tag{MA-RS.56}
\]
The limiting integral and zero sum are absolutely convergent.
The limit is a continuous distribution in $g$ and is independent
of $\chi$ and of the shape of $\phi$ with its fixed endpoints.

Let $R:\mathcal H_a\longrightarrow\mathcal H_{-a}$ denote
physical reflection as in MA-RS-003. Define actual trace-class
operators on the direct sum by
\[
 \mathbb K_N(g)=K_N(g)\oplus R K_N(g^\vee)R.
 \tag{MA-RS.57}
\]
Then
\[
 \lim_{N\to\infty}\operatorname{Tr}\mathbb K_N(g)
   =B_{\overline S}(g)-G(1/2)-G(-1/2)
       +\sum_{|\operatorname{Re}\rho-1/2|>a}
                         m_\rho G(\rho-1/2).
 \tag{MA-RS.58}
\]
There are no boundary zeros at this generic width. These assertions
use neither $A$ nor $M$.

\begin{proof}
\emph{The two operators are defined before the trace.}
The generic-line estimate (MA-RS.2) excludes zeros of $\zeta$
on the line. The pole at one is off it, and the finitely many
removed Euler factors have no zeros there. Thus $A$ and $j$
are smooth on the real frequency axis. All four multipliers
in (MA-RS.54) are smooth and compactly supported, hence bounded.
Since $T_\phi$ is bounded, (MA-RS.54) defines two bounded
operators on all of $L^2$, without any domain assertion for
the uncut expression $J_{a,S}M_\phi$.

The distribution kernel of $T_\phi$ is
$(2\pi)^{-1}\widehat\phi_+(t-u)$, where the hat here is the
unnormalized transform used for $p$ in (MA-RS.53).
On compactly supported smooth frequency vectors, the difference
of the two kernels in (MA-RS.54) is this distribution multiplied
by
\[
 \chi_N(t)\chi_N(u)m(t)\bigl(1-A(t)j(u)\bigr)
  =\chi_N(t)\chi_N(u)m(t)A(t)\bigl(j(t)-j(u)\bigr).
 \tag{MA-RS.59}
\]
Since $\widehat{\phi'}_+(\xi)=-i\xi\widehat\phi_+(\xi)$,
one has $\xi\widehat\phi_+(\xi)=ip(\xi)$ as distributions.
The identity
\[
 j(t)-j(u)=(t-u)\int_0^1j'(u+\theta(t-u))\,d\theta
 \tag{MA-RS.60}
\]
therefore converts (MA-RS.59) exactly into (MA-RS.55).
It also cancels any delta term of $\widehat\phi_+$ at zero.
The integral in (MA-RS.60) is smooth in both variables.
Multiplication by the two compact cutoffs consequently gives
$\mathcal K_N\in C_c^\infty(\mathbb R^2)$.
The resulting integral operator agrees on a dense set with
the already defined bounded difference in (MA-RS.54), hence
agrees everywhere.

\emph{Trace class and the diagonal formula.}
Here is a direct trace-class argument that does not assume a
trace from the diagonal alone. Choose $L$ so that the support
of $\mathcal K_N$ lies in the interior of $[-L,L]^2$.
Use the normalized Fourier basis
$e_r(t)=(2L)^{-1/2}e^{\pi irt/L}$, $r\in\mathbb Z$,
on $[-L,L]$, extended by zero outside that interval.
The smooth kernel is zero near the boundary of the square, so
its periodic Fourier expansion has the form
\[
 \mathcal K_N(t,u)=\sum_{r,s\in\mathbb Z}
              c_{r,s}e_r(t)\overline{e_s(u)},\qquad
 |c_{r,s}|\le
       C_N(1+r^2)^{-1}(1+s^2)^{-1}.
 \tag{MA-RS.61}
\]
The estimate follows by two integrations by parts in each
variable, or by applying $1-\partial_t^2$ and
$1-\partial_u^2$ before taking the Fourier coefficients.
The coefficients are absolutely summable. Accordingly the
associated rank-one operator series converges in trace norm,
and the kernel series converges uniformly on the square.
Its trace is $\sum_r c_{r,r}$, which, by the same absolute
summability and the orthonormality of this basis, equals
$\int\mathcal K_N(t,t)\,dt$.
This proves trace class and the diagonal formula for the
actual operator. No bound uniform in $N$ is needed here.

Finally $p(0)=\int\phi'=-1$ and
$j'(t)=-i\zeta_S'(\sigma+it)/\zeta_S(\sigma+it)^2$.
The diagonal of (MA-RS.55) is therefore
\[
 \mathcal K_N(t,t)=
    -\frac1{2\pi}\chi_N(t)^2G(a+it)
                         \frac{\zeta_S'}{\zeta_S}(\sigma+it).
 \tag{MA-RS.62}
\]
This gives the first line of (MA-RS.56), including its sign
and its Fourier normalization.

\emph{A polynomial bound for the logarithmic derivative.}
We give the needed estimate explicitly; the superpolynomial
upper bound for $1/\zeta_S$ in (MA-RS.3) by itself would not
give it. At large $|t|$ use the finite disk in the proof of
MA-RS-001, centered at $c_t=2+it$ and with a zero-free boundary
of radius $12<R_t<13$. Write $\zeta=H_t B_t$ on that disk,
where $B_t$ is the finite disk Blaschke product over its
$O(\log(2+|t|))$ zeros, with multiplicities.
That proof gives a polynomial upper bound for $|H_t|$
throughout the disk and the polynomial lower bound (MA-RS.5)
on its inner half. At $w=\sigma+it$, Cauchy's estimate on a
fixed-radius circle inside this inner half disk gives a
polynomial bound for $H_t'(w)$. Therefore $H_t'(w)/H_t(w)$
is polynomially bounded.

For a disk Blaschke factor with zero $\rho$, its logarithmic
derivative at $w$ is
\[
 \frac1{w-\rho}+
     \frac{\overline{\rho-c_t}}
          {R_t^2-\overline{\rho-c_t}(w-c_t)}.
 \tag{MA-RS.63}
\]
The unimodular normalizing constant has no derivative.
At $w=\sigma+it$ the second summand is bounded by a fixed
constant, since $|w-c_t|<R_t/4$. The first summand is
$O_\sigma((2+|t|)^2)$ by (MA-RS.2) and the bounded height
range of this disk. Summing over its $O(\log(2+|t|))$
zeros proves a polynomial bound for $B_t'(w)/B_t(w)$.
Adding the two logarithmic derivatives and the bounded
logarithmic derivative of $E_S$ gives constants $C,D$ with
\[
 \left|\frac{\zeta_S'}{\zeta_S}(\sigma+it)\right|
              \le C(2+|t|)^D\qquad(t\in\mathbb R).
 \tag{MA-RS.64}
\]
The compact height range follows from continuity and
nonvanishing on the line.

For arbitrary $g\in C_c^\infty$, repeated integration by
parts shows that $m(t)=G(a+it)$ is Schwartz. Thus its product
with (MA-RS.64) is integrable. Moreover
$|\chi_N(t)|\le\|\chi\|_\infty$ and $\chi_N(t)\to1$
pointwise. Dominated convergence in (MA-RS.62) proves the
second line of (MA-RS.56). The same bound and integration
by parts bound this limiting functional by finitely many
derivative seminorms of $g$ on every fixed compact support.
This proves its continuity, and proves independence of
the frequency cutoff and of the physical cutoff's shape.

\emph{Arithmetic evaluation of the established trace limit.}
Only now rewrite the limiting integral as the contour
integral (MA-RS.29). On a line $\operatorname{Re}s=c>1$
the Euler logarithmic derivative converges absolutely,
and Fourier inversion gives $P_{\overline S}^{+}(g)$
exactly as in (MA-RS.30)--(MA-RS.31).
On the fixed strip $\sigma\le\operatorname{Re}s\le c$,
integration by parts gives
\[
 |G(s-1/2)|\le C_{g,c,\sigma,M}
                    (1+|\operatorname{Im}s|)^{-M}
                  \quad\text{for every integer }M\ge0.
 \tag{MA-RS.65}
\]
Use the selected upper and lower heights of MT-M18
(M18.7)--(M18.8). Its finite-disk minimum-modulus proof
and the fixed-strip estimates give a polynomial bound
for $\zeta_S'/\zeta_S$ on the horizontal segments at
these heights, as checked in MA-RS-003.
Choosing $M$ larger than this polynomial exponent makes
the horizontal integrals tend to zero. On the generic
vertical line absolute convergence follows from
(MA-RS.64); on the line $c$ it follows from the Euler
series and (MA-RS.65).

The counterclockwise rectangle has its right edge directed
upwards and its left edge downwards. The residue of
$G(s-1/2)\zeta_S'(s)/\zeta_S(s)$ is
$m_\rho G(\rho-1/2)$ at a crossed nontrivial zero and
$-G(1/2)$ at the pole at one. Zeros of the removed Euler
factors have real part zero and are not crossed.
Multiplying the resulting left-line integral by
$-(2\pi i)^{-1}$ gives the final line of (MA-RS.56).
The local zero count $O(\log(2+|T|))$ in fixed-height
windows, together with (MA-RS.65), proves absolute
convergence of the zero sum.

Physical reflection is unitary from $\mathcal H_a$ onto
$\mathcal H_{-a}$, so (MA-RS.57) is an actual trace-class
direct sum for each $N$. Apply (MA-RS.56) to $g$ and
$g^\vee$ and add their limits. The identity
$\mathcal M(g^\vee)(z)=G(-z)$ and the multiplicity-preserving
symmetry $\rho\mapsto1-\rho$ turn the second right-hand
zero sum into the sum over $\operatorname{Re}\rho<1-\sigma$.
This proves (MA-RS.58). The symmetry also excludes zeros
on the reflected boundary. No parity of $\chi$ is required:
reflection is applied to the already defined operator
$K_N(g^\vee)$.
\end{proof}

\paragraph{The same finite cutoffs in character sectors.}
For the primitive finite-order Hecke character $\eta$ of
MA-RS-005, choose $\sigma$ in the simultaneous generic set
for $\eta$ and $\overline\eta$. Replace $A(t)$ in
(MA-RS.52)--(MA-RS.55) by
$A_\eta(\sigma+it)$ and $j(t)$ by its reciprocal, and denote
the resulting actual operators by $K_{\eta,N}(g)$.
For arbitrary compact smooth $g$ and $\phi'$ they satisfy
\[
 \begin{aligned}
 \lim_{N\to\infty}\operatorname{Tr}K_{\eta,N}(g)
  &=P^+_{\eta,\overline S}(g)-\delta_\eta G(1/2)
       +\sum_{\substack{\rho\in Z_\eta\\
                         \operatorname{Re}\rho>\sigma}}
                   m_\rho G(\rho-1/2),\\
 \lim_{N\to\infty}\operatorname{Tr}
   \bigl(K_{\eta,N}(g)\oplus
                R K_{\overline\eta,N}(g^\vee)R\bigr)
  &=B_{\eta,\overline S}(g)
         -\delta_\eta\bigl(G(1/2)+G(-1/2)\bigr)\\
  &\qquad +\sum_{\substack{\rho\in Z_\eta\\
                    |\operatorname{Re}\rho-1/2|>a}}
                       m_\rho G(\rho-1/2).
 \end{aligned}
 \tag{MA-RS.66}
\]
Indeed the finite-cutoff boundedness, distributional kernel,
Fourier-series trace-class proof and diagonal evaluation
above use only smooth nonvanishing of $A_\eta$ on the line,
which MA-RS-005 supplies. That proof also supplies its
finite-disk Blaschke factorization, polynomial upper and
zero-free-factor lower bounds, and the distance (MA-RS.46).
Applying exactly (MA-RS.63) to those disk factors gives
a polynomial bound for $A_\eta'/A_\eta$ on the line,
including its bounded finite Euler-factor contribution.
Hence dominated convergence works for every compact smooth
test. The selected-height bounds, Euler derivative and
residues in the last part of MA-RS-005 then give the first
identity of (MA-RS.66), using (MA-RS.65) in place of
Gevrey decay. Its pole has order $\delta_\eta$.
Finally its completed functional equation identifies
$Z_{\overline\eta}$ with $1-Z_\eta$, with multiplicity.
This is precisely the conjugate reflection needed for the
second identity. The root number is a constant and adds
no term. This proves the character extension without
assuming its Hecke GRH.

\paragraph{Compatibility and the precise remaining limit issue.}
When the uncut operator $K_{a,S,\phi}(g)$ of MA-RS-002
is known to be trace class, put
$Q_N=V_a^{-1}M_{\chi_N}V_a$. Its kernel formula and
(MA-RS.55) give the actual equality
\[
           K_N(g)=Q_NK_{a,S,\phi}(g)Q_N.
 \tag{MA-RS.67}
\]
Thus $K_N(g)\to K_{a,S,\phi}(g)$ in trace norm in that
already proved case. To justify this last statement,
$Q_N$ are uniformly bounded self-adjoint operators and
converge strongly to one. The conclusion first holds for
each rank-one operator, since both its defining vectors
converge in norm after applying $Q_N$, then for finite-rank
operators by addition, and then for all trace-class
operators by finite-rank approximation and the uniform
bound on $\|Q_N\|$. It applies to generic-line Gevrey
data by MA-RS-002, or to arbitrary compact smooth data
under the fixed gap of MA-RS-004. The same reasoning
applies in the character cases already covered by
MA-RS-005.

For arbitrary compact smooth data on a merely generic line,
this section proves only the existence and the stated
limit of the actual finite-frequency traces.
It asserts neither trace-norm nor operator convergence
of $K_N(g)$ in that larger scope. The regularization has
changed by inserting frequency cutoffs into the two
bounded compositions. It does not identify their limit
with Masullo's physical cutoff, an orthogonal quotient
projection, or all off-diagonal maps in Meyer's common
space. The polar and exterior-divisor corrections in
(MA-RS.58) and (MA-RS.66) remain part of the formula.

\paragraph{Sources and geometric placement.}
The generic set, disk bounds and selected heights are
the unconditional inputs of MA-RS-001--003, with their
exact inherited source NT-Z10 and MT-M18
(M18.5)--(M18.9). The character extension uses the
Hecke inputs and exact conventions of MA-RS-005,
including Ramakrishnan and Valenza (1999),
\emph{Fourier Analysis on Number Fields}, \S7.4,
Theorem~7-19, p.~279, and \S7.7, Theorem~7-28, p.~289,
through the already checked inherited consequences.
The bounded-composition, compact-kernel, diagonal and
limit arguments are proved above. Meyer's original
wide-strip diagonal trace is located at Meyer (2005),
Lemmas~5.3--5.6 and Theorem~5.8, author pp.~39--44,
as checked in MA-MEY-001--003; no new narrow-width
claim is attributed to that source.
In a character sector, the prime coefficients retain
the Frobenius-trace interpretation of MA-RS-005 and
MA-CH-001: SGA~1, Expos\'e~V, \S7, and Deligne,
\emph{Valeurs de fonctions L}, Definition~6.1 and
\S\S6.2--6.3. The trace is the trace of the actual
regularized analytic operators constructed here;
this interpretation does not supply positivity or
zero location.
